% Created 2026-06-12 vie 10:03
% Intended LaTeX compiler: pdflatex
\documentclass{article}
\usepackage[utf8]{inputenc}
\usepackage[T1]{fontenc}
\usepackage{graphicx}
\usepackage{longtable}
\usepackage{wrapfig}
\usepackage{rotating}
\usepackage[normalem]{ulem}
\usepackage{amsmath}
\usepackage{amssymb}
\usepackage{capt-of}
\usepackage{hyperref}

\usepackage{fancyhdr}
\usepackage[top=25mm, left=25mm, right=25mm]{geometry}
\usepackage{longtable}
\fancyhead[R]{}
\setlength\headheight{43.0pt}
\usepackage{tabularx}
\usepackage[T1]{fontenc}
\usepackage[utf8]{inputenc}
\usepackage[spanish]{babel}
\usepackage[backend=biber]{biblatex}
\addbibresource{../bibliography/bibliography.bib}
\author{Danilo UnapuCha}
\date{2026-06-10}
\title{Taller — Máquina IAS / Von Neumann\\\medskip
\large ICCD332 Arquitectura de Computadores}
\hypersetup{
 pdfauthor={Danilo UnapuCha},
 pdftitle={Taller — Máquina IAS / Von Neumann},
 pdfkeywords={},
 pdfsubject={},
 pdfcreator={Emacs 27.1 (Org mode 9.7.5)}, 
 pdflang={Spanish}}
\begin{document}

\maketitle
\fancyhead[C]{\includegraphics[scale=0.05]{../images/logoEPN.jpg}\\
ESCUELA POLITECNICA NACIONAL\\FACULTAD DE INGENIERIA DE SISTEMAS\\
ARQUITECTURA DE COMPUTADORES}
\thispagestyle{fancy}
\section{Objetivos}
\label{sec:org5a45bc4}

\begin{itemize}
\item Comprender el ciclo de captacion, decodificacion y ejecucion de la
maquina IAS (arquitectura Von Neumann).
\item Identificar y aplicar el conjunto de instrucciones IAS para traducir
codigo maquina a lenguaje ensamblador y a notacion RTL.
\item Simular mentalmente la ejecucion de un programa en la maquina IAS
para determinar el estado final de los registros.
\end{itemize}
\section{Instrucciones}
\label{sec:orgcb8bb24}

Complete las solicitudes del taller y suba el archivo \texttt{.org} y \texttt{.pdf}
al aula virtual. Use las combinaciones de teclado para trabajar en Emacs.
\section{Recursos}
\label{sec:org65d76f7}

Para este taller consulte el libro \citetitle{stallings2006esp}
disponible en \href{https://epnecuador-my.sharepoint.com/:f:/g/personal/lenin\_falconi\_epn\_edu\_ec/EgjH2RoedD5NuswqpOt8ExsB\_DE052v9Rlrg0QpEtbimDg?e=WoqexR}{EPN-SHAREPOINT (Clic aqui)} y use su WSL con Emacs para editar el archivo.
\subsection{Observaciones y Recomendaciones}
\label{sec:org9f0df6b}
\begin{itemize}
\item Tenga en cuenta la ruta a las imagenes y a la bibliografia. Se le
recomienda que realice una clonacion o un fork del repositorio de la
clase a fin de que no tenga problemas generando el documento.
\item Para la generacion correcta de archivos la configuracion basica debe
estar operativa y funcional (es decir, Emacs, \LaTeX).
\end{itemize}
\section{Taller: Maquina de Von Neumann}
\label{sec:orgb32cda4}

La maquina de Von Neumann realiza operaciones en un ciclo repetitivo de
\textbf{captacion}, \textbf{decodificacion} y \textbf{ejecucion}. Para esto, el computador
IAS usa una memoria de 40 bits. Cuando la memoria se usa
numericamente, el dato representa un número en complemento a 2; donde
el bit inicial corresponde al signo. Cuando la memoria se usa para el
registro de instrucciones, se divide en dos partes de 20bits cada una. Los 8
primeros bits de cada parte corresponden al \emph{opcode} y los restantes
12 bits de cada parte a la \emph{direccion (operando)} desde la que se debe
leer o a la que se debe escribir.
\subsection{Conjunto de Instrucciones IAS}
\label{sec:orgba7d5a6}

Complete la Tabla \ref{tab:tablaISA-IAS}  de instrucciones del computador IAS
\autocite[p.47]{stallings2006esp} usando la notacion RTL\footnote{Register Transfer Language o Lenguaje de Transferencia de Registros.}. Se suministra
la Tabla en codigo \LaTeX{} para tal efecto.

Para editar la Tabla \ref{tab:tablaISA-IAS} ubiquese en
el bloque \texttt{\#+begin\_export latex} y ejecute el comando \texttt{C-c'}. Esto
dividira el buffer en dos ventanas. El cursor salta a la nueva ventana
y le permite editar la tabla. Dentro de la nueva tabla el modo
predominante cambiara automaticamente a \LaTeX, lo que le permite
hacer cambios con la asistencia del IDE. Para salir y aceptar los
cambios ejecute de nuevo \texttt{C-c'}.

\begin{itemize}
\item \texttt{M} equivale a \emph{Alt}
\item \texttt{x} es la tecla \emph{x}
\item \texttt{RET} es presionar \emph{Enter}
\end{itemize}

Una vez activado el modo \LaTeX{} el siguiente codigo recibe colores
sobre las palabras clave.

\begin{table}[htbp]
  \caption{Instrucciones Maquina IAS}
  \label{tab:tablaISA-IAS}
  \centering
  \begin{tabular}{|cccc|}
    \hline
    Opcode   & Opcode Hex &  Simbolo          & RTL \\ \hline
    00001010 & 0x0A       &  LOAD MQ          & $[AC]\leftarrow[MQ]$\\
    00001001 & 0x09       &  LOAD MQ, M(X)    & $[MQ]\leftarrow[X]$\\
    00100001 & 0x21       &  STOR M(X)        & $[X]\leftarrow[AC]$\\
    00000001 & 0x01       &  LOAD M(X)        & $[AC]\leftarrow[X]$\\
    00000011 & 0x03       &  LOAD |M(X)|      & $[AC]\leftarrow|[X]|$\\
    00000100 & 0x04       &  LOAD -|M(X)|     & $[AC]\leftarrow-|[X]|$\\
    \hline
  \end{tabular}
\end{table}

\newpage
\subsection{Ejercicio}
\label{sec:orgd928b14}

En la maquina IAS, las instrucciones se dividen en dos segmentos:
i) izquierdo, desde el bit 0 a 19, y ii) derecho, desde el bit 20 al 39. Primero
se ejecuta el lado izquierdo y luego el derecho. El
contador de programa inicia en la posicion \([PC] \leftarrow 300\). El conjunto de
instrucciones del computador o ISA\footnote{Instruction Set Architecture: conjunto de instrucciones que la maquina IAS puede interpretar y ejecutar.} esta definido en la Tabla
\ref{tab:tablaISA-IAS}. A fin de terminar el programa, se agrega la
instruccion \texttt{0x99} que lleva al computador al estado de HALT. Conteste las
siguientes preguntas considerando la Tabla \ref{tab:org4511338} como el
mapa de memoria de la maquina y obtenga el programa en lenguaje ensamblador y en
RTL.

\begin{table}[htbp]
\caption{\label{tab:org4511338}Mapa de Memoria}
\centering
\begin{tabular}{rrrrr}
\hline
Direccion & \(Opcode_1\) & \(X_1\) & \(Opcode_2\) & \(X_2\)\\
\hline
0x300 & 01 & 940 & 06 & 941\\
0x301 & 21 & 940 & 14 & 000\\
0x302 & 21 & 941 & 02 & 940\\
0x303 & 21 & 940 & 99 & 000\\
\ldots{} & \ldots{} & \ldots{} & \ldots{} & \ldots{}\\
0x940 & 00 & 000 & 00 & 005\\
0x941 & 00 & 000 & 00 & 002\\
\hline
\end{tabular}
\end{table}
\subsubsection{Preguntas}
\label{sec:org76be429}
Procedmiento:
Para la direccion \(\text{0x300}\)
\begin{itemize}
\item \([PC] \leftarrow [IR]\)
\item \([IR] \leftarrow \text{0x300}\)
\item \([IR] \leftarrow \text{01940}\)
\item \([AC] \leftarrow [940]\)
\item \([AC] \leftarrow \text{5}\)

\item \([PC] \leftarrow [IR]\)
\item \([IR] \leftarrow \text{0x300}\)
\item \([IR] \leftarrow \text{06941}\)
\item \([AC] \leftarrow [AC]-[941]\)
\item \([AC] \leftarrow 5-2\)
\item \([AC] \leftarrow \text{3}\)
\end{itemize}

Para la direccion \(\text{0x301}\)

\begin{itemize}
\item \([PC] \leftarrow [IR]\)
\item \([IR] \leftarrow \text{0x301}\)
\item \([IR] \leftarrow \text{21940}\)
\item \([940] \leftarrow [AC]\)
\item \([940] \leftarrow \text{3}\)

\item \([PC] \leftarrow [IR]\)
\item \([IR] \leftarrow \text{0x301}\)
\item \([IR] \leftarrow \text{14000}\)
\item \([AC] \leftarrow [AC] * 2\)
\item \([AC] \leftarrow \text{6}\)
\end{itemize}

Para la direccion \(\text{0x302}\)

\begin{itemize}
\item \([PC] \leftarrow [IR]\)
\item \([IR] \leftarrow \text{0x302}\)
\item \([IR] \leftarrow \text{21941}\)
\item \([941] \leftarrow [AC]\)
\item \([941] \leftarrow \text{6}\)

\item \([PC] \leftarrow [IR]\)
\item \([IR] \leftarrow \text{0x302}\)
\item \([IR] \leftarrow \text{02940}\)
\item \([AC] \leftarrow -[940]\)
\item \([AC] \leftarrow \text{-3}\)
\end{itemize}

Para la direccion \(\text{0x303}\)

\begin{itemize}
\item \([PC] \leftarrow [IR]\)
\item \([IR] \leftarrow \text{0x303}\)
\item \([IR] \leftarrow \text{21940}\)
\item \([940] \leftarrow [AC]\)
\item \([940] \leftarrow \text{-3}\)

\item \([PC] \leftarrow [IR]\)
\item \([IR] \leftarrow \text{0x303}\)
\item \([IR] \leftarrow \text{99000}\)
\item \(\text{HALT}\)
\end{itemize}

\begin{enumerate}
\item ¿Que resultado se tiene al final en el registro del acumulador?  
Respuesta:

\([AC] \leftarrow -3\)
\end{enumerate}


\begin{enumerate}
\item ¿Se sobrescribe algún registro como resultado de la ejecucion? Si
es verdadero, indique que registro y con que valor.
Respuesta:
\begin{itemize}
\item \([940] \leftarrow -3\)
\item \([941] \leftarrow 6\)
\item \([AC] \leftarrow -3\)
\end{itemize}

\item ¿Cual es el valor del bit de Carry? ¿Que operacion modifica el bit?
Respuesta:
\begin{itemize}
\item \(Carry = 1\)

\item La operacion que modifica el bit es SUB M(941)
Aqui se hacemos que la ALU haga la siguiente operacion \([AC]
     \leftarrow [AC] + Comp2([M(941)])\)
Donde tenemos que el valor de \([AC]\) es 5 o (\(00000101\)) y el
valor de \(M(941)\) es 2 o (\(00000010\)), a este mismo valor se le
hace el complemento a 2 dando como resultado (\(11111110\)) y ahi
si se va a poder hacer la operacion solicitada por la direccion
de memoria siendo esta la resta que se veria asi
  00000101
\begin{itemize}
\item 11111110

\noindent\rule{\textwidth}{0.5pt}
100000011
\end{itemize}
Aqui ya podemos observar que existe el carry de uno al salirse de
nuestra arquitectura de 8 bits
\end{itemize}
\end{enumerate}
\subsubsection{Codigo RTL}
\label{sec:org773b3cf}

Escriba en notacion de transferencia de registros el programa que
ejecuta el computador IAS.

\begin{itemize}
\item \([AC] \leftarrow [940]\)
\item \([AC] \leftarrow [AC] - [941]\)
\item \([940] \leftarrow [AC]\)
\item \([AC] \leftarrow [AC] \times 2\)
\item \([941] \leftarrow [AC]\)
\item \([AC] \leftarrow -[940]\)
\item \([940] \leftarrow [AC]\)
\item \(HALT\)
\end{itemize}
\subsubsection{Codigo Assembler}
\label{sec:orgad573f6}

Escriba el codigo Assembler del programa que ejecuta el computador IAS.

\begin{verbatim}
LOAD M(940)
SUB M(941)
STOR M(940)
SHL AC,1
STOR M(941)
LOAD M(940)
NEG
STOR M(940)
HALT
\end{verbatim}
\section{Equipo de Trabajo}
\label{sec:orgf74dcdf}

Complete la informacion de los integrantes del grupo.

\begin{center}
\begin{tabular}{lll}
\hline
Nombre & Email & Rol\\
\hline
Danilo Unapucha & danilo.unapucha@epn.edu.ec & Lider\\
Michael Sornoza & michael.sornoza@epn.edu.ec & Colaborador\\
Sebastian Arevalo & martin.arevalo@epn.edu.ec & Colaborador\\
\hline
\end{tabular}
\end{center}
\section{Verificacion de Entregables [100\%]}
\label{sec:orgf8f8c97}

Ejecute \texttt{C-c C-c} sobre los items de tarea según se hayan cumplido o
no. Si un item no pudo realizarse, anote en la seccion de problemas
las razones.

\begin{itemize}
\item[{$\boxtimes$}] Tabla de instrucciones IAS completada.
\item[{$\boxtimes$}] Preguntas del ejercicio respondidas.
\item[{$\boxtimes$}] Codigo RTL completo.
\item[{$\boxtimes$}] Codigo Assembler completo.
\item[{$\boxtimes$}] Revision de ortografia con \texttt{ispell-buffer}.
\item[{$\boxtimes$}] Exportacion a PDF con \texttt{M-x org-latex-export-to-pdf}.
\item[{$\boxtimes$}] Subir al aula virtual archivo \texttt{.org} y \texttt{.pdf}.
\end{itemize}
\subsection{Problemas con la tarea / logisticos}
\label{sec:org41116d4}

\begin{itemize}
\item Tarea \(X_1\) no pudo completarse debido a:
\item Tarea \(X_2\) no pudo completarse debido a:
\end{itemize}

\printbibliography
\end{document}
